Tá istá sekvencia bude pravdepodobne načítaná viackrát, avšak s odlišnými chybami, keďže nanopórové sekvenovanie je chybovejšie.
Okrem toho sa sekvencia môže načítať v opačnej orientácii, pretože molekula vstupuje do nanopóru z rôznych koncov.
Sekvenciu preto možno získať v smere 3$'$\textrightarrow5$'$ alebo 5$'$\textrightarrow3$'$.

V dôsledku biosyntézy vzniká aj komplementárna sekvencia DNA, v ktorej sú bázy nahradené podľa pravidiel párovania (A~\textrightarrow~T, T~\textrightarrow~A, G~\textrightarrow~C, C~\textrightarrow~G).
Aj komplementárna sekvencia môže byť čítaná v oboch smeroch.
Pre každú pôvodnú sekvenciu sa teda môžu vyskytnúť až štyri varianty: pôvodná sekvencia v smere 3$'$\textrightarrow5$'$ a 5$'$\textrightarrow3$'$, a jej komplement v smere 3$'$\textrightarrow5$'$ a 5$'$\textrightarrow3$'$.
